\newlang{\SDD}{SDD}
\newlang{\DSOS}{DSOS}
\newlang{\SDSOS}{SDSOS}
\newlang{\SOS}{SOS}


In this paper, they provide heuristics to show that a polynomial 

\[p(x)=\sum_i c_{\alpha_1,\dots,\alpha_n}x_1^{\alpha_1}\dots x_n^{\alpha_n} \geq 0\]

And to use such approximations to perform polynomial optimization (i.e., minimization or maximization of a polynomial function, given the coefficients $c_{\alpha_1,\dots,\alpha_n}$ understand if the polynomial is always positive?).

It is important to notice that a \textit{form} or \textit{homogeneous} polynomial is a polynomial where all monomials have the same degree.

To begin with, they provide different approximations to show such a fact.

A sufficient condition is to write $p(x) = \sum q_i(x)^2$ as a SOS (sum of squares), with the degree of $p$ to be $2d$.

This can also be seen as $p(x) = z^T(x,d) Q z(x, d)$, with $z(x,d)$ the vector of all monomials in $x$ of degree $d$, and have $Q \succeq 0$.

Of course, since the condition $Q \succeq 0$ is expensive, we are looking for approximations, hence we will rather use $Q \in \DD$ or $Q \in \SDD$ (i.e., scaled diagonally dominant, namely if there exists a diagonal matrix $D$ with nonnegative entries s.t. $D Q D\in \DD$).

So, we define the three cones of nonnegative polynomials with $n$ variables and degree $d$ as:

\begin{align}
	\DSOS_{n,d} \subseteq \SDSOS_{n,d} \subseteq \SOS_{n,d} \subseteq \PSD_{n,d}
\end{align}

And the goal is to provide methods to determine how to find these matrices $Q$.

In section 3.1 \hubi{which is the one that Monaldo told me to study}, 

This theorem shows why these approximations are useful:

\begin{theorem}
\label{thm:poly_optim_approx_is_lp}
	For a fixed degree $d$, optimization over $\DSOS_{n,d}$ or $\SDSOS_{n,d}$ can be done using LP or SOCP (Second Order Cone Programming).
\end{theorem}

So, when we use another condition different than $X\succeq 0$, the optimization problem becomes suddenly easy.

In Section 3, they develop an example which is useful to understand what kind of approximation they use: suppose you are given the problem $\SOS^*$

\begin{align}
	\SOS^* := & \min_{X\in S_n} & X \cdot C & \\
	& \text{ s.t. } & A_i \cdot X = b_i & & i=1,\dots,m \\
	& & X\succeq 0 & \\
\end{align}

which is clearly an SDP, and you substitute the condition of positive semidefiniteness with $X\in\DD(U)$, where 
\begin{definition}
	$\DD(U) := \{M\in S_n: \exists Q\in\DD, M=U^T Q U\}$.
\end{definition}

And call this new optimization problem $\DSOS_k$.

They say it is clear that optimizing over $\DD(U)$ is an LP (fair enough, as long as you turn all the absolute values of the $\DD$ condition into an LP, which I am pretty sure you can do easily).

Now, the algorithm they propose is: start with $U_1=I$ the identity matrix, then you compute $X_1$ the optimizer, and you proceed iteratively computing $U_{k+1} = chol(X_k)$ the Cholesky decomposition of $X_k=U_{k+1}^T U_{k+1}$.

It can be proved that during this iterative procedure, the value of matrices $X_k$ can only improve (we are kindof refining the search cone $\DD(U_k)$), so the sequence $\{\DSOS_k\} \geq \SOS^*$ is decreasing, and converging to a value $\DSOS^*\geq \SOS^*$. They claim they could not prove is how much worse than $\SOS^*$ is $\DSOS^*$, but experimentally it works well.

